\section{Conclusion}

In this work, we studied transferability between understanding and generation in UMMs. Through controlled experiments, we showed that capabilities learned in one modality can transfer to the other, and that the strength of this transfer strongly depends on the architectural design of the model. Building on this observation, we demonstrated that transferability can be exploited as a practical training strategy. Instead of directly fine-tuning generation models for specific capabilities, introducing these capabilities through understanding tasks allows the knowledge to transfer into generation while preserving overall generative quality. Experiments across counting, spatial relation reasoning, and OCR tasks show that this approach can effectively improve generative capabilities without degrading general multimodal performance. We hope our findings provide new insights into the interaction between understanding and generation and open new directions for post-training in unified multimodal models.
